\documentclass[11pt,a4paper]{article}
\usepackage[margin=25mm]{geometry}
\usepackage{amsmath,amssymb,amsthm,hyperref}
\setlength{\emergencystretch}{2em}
\newtheorem{theorem}{Theorem}
\newtheorem{corollary}[theorem]{Corollary}
\title{A split admitting a rank and degree extension}
\author{Human Compression project}\date{}
\begin{document}\maketitle
A perfect split is internal when its label is neither endpoint label.
Its crossing mark has the form
\[
                    q=(a\ b)\alpha\beta,
\]
where the prefix and suffix are supported and transitive on disjoint
sets $A,B$, the crossing points belong to their respective sets,
and the internal factors are supported off $a,b$ on the corresponding
sides. Call the split \emph{sufficient} if $q$ is pure, or if one
internal factor is trivial and both actual attaching groups contain
an odd permutation. When $\alpha=1$ these groups are the prefix
$L$ and the group $J$ generated by $q$ and the suffix. Their supports
meet at $a$. This definition retains the exact hypotheses of the
acquired pure-split and odd-attachment theorems.

The acquired components include the complete parity proof for the
remaining middle, its exact finite boundaries and their distinct reviews.
We now close the induction on the actual degree.

\begin{theorem}\label{sufficient}
Let a fractured string C-marking of the natural symmetric group
$S_n$, with $n\ge5$, have rank $r$ satisfying $2r>n+2$.
It has an internal sufficient split.
For $n\ge7$, the fracture hypothesis follows from the same rank
inequality and need not be imposed separately.
\end{theorem}
\begin{proof}
We use strong induction on the actual degree. The acquired orbit
forest gives $r\le n-1$, so degrees at most four cannot meet the
strict inequality. In degrees five and six, that inequality forces
rank $n-1$. The acquired complete extremal classification gives
the ordinary adjacent-transposition simplex, apart from the outer
degree-six action. The latter has a transitive maximal omission
and is excluded by fracture. Every internal simplex mark is a pure
perfect split. These are the complete induction bases.

Suppose $n\ge7$ and the statement is known at every smaller degree.
Apply the acquired conditional reduction to even anchored factors.
Its complete hypotheses are exactly the present natural symmetric
marking, fracture, strict rank inequality and smaller-degree
sufficient-split induction. It uses the already proved trivial-factor
theorem, actual-wing theorem, even-attachment bound and one-point
transport, so no split-existence or quotient-intersection assumption
is being added here. Its conclusion is a sufficient split, or the
following actual residual, after reversal:
\[
 \begin{array}{c|c|c}
 \text{factor}&\text{degree}&\text{rank}\\\hline
 P\text{, natural symmetric}&m\text{ even}&m/2+1\\
 E\text{, even}&e\in\{6,8,10\}&e/2.
 \end{array}
\]
The label intervals are consecutive and disjoint; both factors are
fractured and transitive on supports meeting only at their common
endpoint anchor. Their full degree and rank satisfy $n=m+e-1$
and $2r=n+3$.

The acquired small-even theorem applies to $E$. Its external group
is the actual odd transitive group $P$ on the other support, meeting
it only at the anchor, within the original C-marking. These are
precisely its attachment and supported-kernel hypotheses. Thus
$e=6$ and $E$ is the primitive degree-six rank-three $A_5$ factor.
The complete first-cut and endpoint arguments, followed by the
paired-middle parity theorem, supply a sufficient split in exactly
this critical configuration. Their uses of induction are exclusively
in actual symmetric factors of degree strictly less than the present
$n$, with the same strict inequality. Their separate rank induction
for paired middles uses only anchor deletion and finite boundaries,
not the theorem now being proved. This completes strong induction.

Finally the acquired initial-fracture theorem applies to every
natural symmetric marking with $n\ge7$ and $2r>n+2$, with its
complete exceptional-rank inventory retained. It supplies the last
assertion without a new finite premise.
\end{proof}

\begin{corollary}\label{extension}
For every natural symmetric string C-marking of degree $n\ge5$
and rank $r$ with $2r>n+2$, there is a natural symmetric string
C-marking of degree $n+1$ and rank $r+1$ obtained by subdividing
an internal perfect crossing. At degree six the marking may first
be replaced by an abstractly isomorphic natural representative.
\end{corollary}
\begin{proof}
For $n\ge7$, Theorem~\ref{sufficient} supplies the split. If it
is pure, use the acquired pure-split theorem. Otherwise reverse
labels if needed to make $\alpha=1$. At a new point $c$, replace
\[
                  q=(a\ b)\beta
       \quad\hbox{by}\quad (a\ c),\ (c\ b)\beta.
\]
The acquired odd-attachment theorem verifies both complete leaf
kernel families from the original IP and the two odd attaching
groups. Its exact split criterion then gives the full intersection
property and string relations. The new group contains the old
$S_n$ fixing $c$, since conjugating the second new mark by the
first recovers $q$. It also contains the pure transposition $(a\ c)$,
so it is $S_{n+1}$. Rank and degree both increase by one.

At degrees five and six the full extremal classification again
reduces the abstract marking to the simplex, whose subdivision is
the next simplex. The outer action of $S_6$ is abstractly isomorphic
to that marking, though it was correctly excluded from the fractured
natural-action theorem. This accounts for the distinction between
actual permutation actions and abstract string C-markings.
\end{proof}

\paragraph{Attribution and unproved counting step.}
The source is Cameron--Fernandes--Leemans, \emph{The number of
string C-groups of high rank}, arXiv:2212.12723v2. The actual
subdivision is their rank and degree operation. All rank
classifications and finite lists remain attributed in the acquired
components. The present proof combines the repaired restriction
criterion, odd-attachment lemma and the completed critical route.
It claims existence of a valid extension. It does not yet prove
an injection or bijection on markings modulo isomorphism and
duality: an intrinsic choice, a valid contraction, and compatibility
with reversal remain to be established. No whole-original credit,
new conjecture, Lean theorem or compression measurement is assigned.
\end{document}
